\subsubsection{DatasetListPage}
\textbf{Componente padre}
\begin{itemize}
    \item StandardPage.
\end{itemize}
\textbf{Modulo}
\begin{itemize}
    \item DatasetModule, modulo funzionale dell'applicazione che raccoglie e organizza tutti i componenti, servizi e risorse legati alla creazione, visualizzazione e gestione di dataset.
\end{itemize}
\textbf{Elementi chiave}
\begin{itemize}
    \item \textbf{SearchDataset}: barra per la ricerca dei dataset, realizzata tramite sottocomponente \textbf{SearchBar};
    \item \textbf{DatasetListContainer}: lista scorrevole dei dataset salvati, realizzata tramite sottocomponente \textbf{DatasetListView};
    \item \textbf{CreateDatasetButton}: pulsante che permette la creazione di un dataset temporaneo;
    \item \textbf{CreateDatasetFromJSONButton}: pulsante che permette di richiedere la creazione di un dataset a partire da un file JSON, la creazione viene poi realizzata tramite \textbf{JsonFileUpload};
    \item \textbf{StatusMessage}: popup di conferma di eliminazione di un dataset o di sovrascrittura del dataset caricato (se modificato), realizzato tramite sottocomponente \textbf{ConfirmComponent};
    \item \textbf{ErrorMessage}: messaggio di errore visualizzato in seguito a un errore nell'ottenimento o nel salvataggio dei dati, realizzato tramite sottocomponente \textbf{MessageComponent}.
\end{itemize}
\textbf{Servizio}
\begin{itemize}
    \item \textbf{DatasetService}: servizio iniettato tramite `@Injectable`, fornisce le funzionalità principali per la gestione della lista di dataset.
\end{itemize}

\subsubsection{DatasetListView}
\textbf{Componente padre}
\begin{itemize}
    \item DatasetListPage.
\end{itemize}
\textbf{Elementi chiave}
\begin{itemize}
    \item \textbf{StatusMessage}: paragrafo contenente il messaggio che viene visualizzato dall'utente nel caso non ci siano dataset salvati nel sistema;
    \item \textbf{DatasetList}: lista scorrevole dei dataset salvati nel sistema, permette la creazione di un sottocomponente \textbf{DatasetElement} per ogni dataset salvato (tramite la direttiva \texttt{*ngFor}).
\end{itemize}
\textbf{Input}
\begin{itemize}
    \item data sets: lista contenente tutti i dataset salvati e/o ricercati nel sistema.
\end{itemize}
\textbf{Output}
\begin{itemize}
    \item rename Dataset: evento emesso quando riceve da \textbf{DatasetElement} la richiesta di rinomina di un dataset;
    \item copy Dataset: evento emesso quando riceve da \textbf{DatasetElement} la richiesta di copia di un dataset;
    \item delete Dataset: evento emesso quando riceve da \textbf{DatasetElement} la richiesta di eliminazione di un dataset;
    \item load Dataset: evento emesso quando riceve da \textbf{DatasetElement} la richiesta di caricare un dataset.
\end{itemize}

\subsubsection{DatasetElement}
\textbf{Componente padre}
\begin{itemize}
    \item DatasetListView.
\end{itemize}
\textbf{Elementi chiave}
\begin{itemize}
    \item \textbf{NameText}: sezione testuale contenente il nome del dataset;
    \item \textbf{DatasetDate}: sezione testuale contenente la data dell'ultima modifica del dataset;
    \item \textbf{RenameButton}: pulsante per richiedere la rinominazione del dataset, questa viene poi realizzata tramite il sottocomponente \textbf{DatasetNameDialog};
    \item \textbf{CopyButton}: pulsante per richiedere la copia del dataset;
    \item \textbf{DeleteButton}: pulsante per richiedere l'eliminazione del dataset, questa viene poi confermata o annullata tramite il sottocomponente \textbf{ConfirmComponent};
    \item \textbf{LoadButton}: pulsante per richiedere il caricamento del dataset, se il dataset attualmente caricato presenta modifiche viene richiesta la conferma o l'annullamento della sovrascrittura tramite il sottocomponente \textbf{ConfirmComponent}.
\end{itemize}
\textbf{Input}
\begin{itemize}
    \item dataset: dataset associato al componente.
\end{itemize}
\textbf{Output}
\begin{itemize}
    \item rename: evento emesso quando viene premuto il pulsante per richiedere la rinominazione del dataset;
    \item copy: evento emesso quando viene premuto il pulsante per richiedere la copia del dataset;
    \item delete: evento emesso quando viene premuto il pulsante per richiedere l'eliminazione del dataset;
    \item load: evento emesso quando viene premuto il pulsante per richiedere il caricamento del dataset.
\end{itemize}

\subsubsection{DatasetNameDialog}
\textbf{Componente padre}
\begin{itemize}
    \item DatasetListPage;
    \item DatasetContentPage.
\end{itemize}
\textbf{Elementi chiave}
\begin{itemize}
    \item \textbf{NewNameLabel}: etichetta associata al campo di inserimento del nuovo nome;
    \item \textbf{NewNameInput}: campo di inserimento del nuovo nome da assegnare al dataset;
    \item \textbf{ConfirmRenameButton}: pulsante che permette di confermare la richiesta di assegnazione o modifica del nome del dataset;
    \item \textbf{CancelRenameButton}: pulsante che permette di annullare la richiesta di assegnazione o modifica del nome del dataset.
\end{itemize}
\textbf{Input}
\begin{itemize}
    \item currentName: gli viene assegnato il nome del dataset nel caso non sia temporaneo.
\end{itemize}
\textbf{Output}
\begin{itemize}
    \item confirmRename: evento emesso quando viene premuto il pulsante per confermare l'assegnazione o modifica del nome del dataset;
   %\item cancelRename: evento emesso quando viene premuto il pulsante per annullare l'assegnazione o modifica del nome del dataset. 
\end{itemize}
